\documentclass[12pt]{article}
\usepackage{fullpage}
\usepackage{amsmath,amssymb,amsthm}
\usepackage{hyperref}

\newtheorem{theorem}{Theorem}
\newtheorem{lemma}[theorem]{Lemma}
\newtheorem{proposition}[theorem]{Proposition}
\newtheorem{remark}[theorem]{Remark}

\title{Solution to Problem 2:\\Whittaker Functions and Rankin--Selberg Integrals}
\author{}
\date{April 14, 2026}

\begin{document}
\maketitle

\section*{Answer}

\textbf{Yes.} There exists $W \in \mathcal{W}(\Pi, \psi^{-1})$ with the stated property.

\section*{Proof}

\subsection*{Setup and notation}

Let $F$ be a non-archimedean local field with ring of integers $\mathfrak{o}$ and uniformizer $\varpi$.  Let $q = |\mathfrak{o}/(\varpi)|$.  Let $\Pi$ be a generic irreducible admissible representation of $G' = \mathrm{GL}_{n+1}(F)$ and $\pi$ a generic irreducible admissible representation of $G = \mathrm{GL}_n(F)$, both over $\mathbb{C}$.

Let $\mathfrak{q}$ be the conductor ideal of $\pi$ and $Q \in F^\times$ a generator of $\mathfrak{q}^{-1}$.  Set
\[
  u_Q = I_{n+1} + Q\,E_{n,n+1} \in G'.
\]

\subsection*{Local Rankin--Selberg theory}

The local Rankin--Selberg integral for $\mathrm{GL}_{n+1} \times \mathrm{GL}_n$ was introduced by Jacquet--Piatetski-Shapiro--Shalika (JPSS).  For $W \in \mathcal{W}(\Pi,\psi^{-1})$ and $V \in \mathcal{W}(\pi,\psi)$, the integral
\[
  \Psi(s;W,V) := \int_{N_n\backslash G}
    W\!\bigl(\mathrm{diag}(g,1)\,u_Q\bigr)\,V(g)\,|\det g|^{s-\frac12}\,dg
\]
converges absolutely for $\mathrm{Re}(s) \gg 0$, admits meromorphic continuation to all of $\mathbb{C}$, and lies in
\[
  \mathbb{C}[q^s, q^{-s}]\cdot L(s,\Pi\times\pi),
\]
where $L(s,\Pi\times\pi)$ is the Rankin--Selberg local $L$-factor.

\subsection*{Essential vectors}

\begin{theorem}[Existence of essential vector]
  \label{thm:essential}
  There exists $W_0 \in \mathcal{W}(\Pi,\psi^{-1})$ such that for every $V_0 \in \mathcal{W}(\pi,\psi)$ equal to the normalized newvector of $\pi$ (the unique-up-to-scalar $K_n(\mathfrak{q})$-fixed vector, normalized by $V_0(I_n)=1$), the identity
  \[
    \Psi(s;\,W_0,\,V_0) = L(s,\Pi\times\pi)
  \]
  holds for all $s \in \mathbb{C}$.
\end{theorem}

\begin{proof}
This is essentially the content of the theory of \emph{essential vectors} (also called \emph{test vectors} or \emph{newvectors}) for the Rankin--Selberg convolution.  The construction proceeds in two steps.

\medskip
\noindent\textbf{Step 1: Existence for unramified $\pi$.}
If $\pi$ is unramified ($\mathfrak{q} = \mathfrak{o}$, so $Q = \varpi^0 = 1$ and $u_Q = I_{n+1}$), then $V_0 = V_0^{\mathrm{sph}}$ is the spherical vector and $W_0 = W_0^{\mathrm{sph}}$ is the spherical Whittaker function.  The Casselman--Shalika formula gives
\[
  W_0^{\mathrm{sph}}\!\bigl(\mathrm{diag}(\varpi^{a_1},\ldots,\varpi^{a_{n+1}})\bigr)
  = \delta_B^{1/2}\!\bigl(\mathrm{diag}(\varpi^{a_i})\bigr)\cdot s_\lambda(\alpha_1,\ldots,\alpha_{n+1}),
\]
where $s_\lambda$ is the Schur polynomial in the Satake parameters of $\Pi$ and $\lambda=(a_1,\ldots,a_{n+1})$.  The JPSS computation then yields $\Psi(s;W_0^{\mathrm{sph}}, V_0^{\mathrm{sph}}) = L(s,\Pi\times\pi)$.

\medskip
\noindent\textbf{Step 2: Ramified $\pi$, the $u_Q$-twist.}
For general $\pi$ with conductor $\mathfrak{q} = (\varpi^a)$, the role of $u_Q$ is to ``undo'' the ramification: the matrix $u_Q = I_{n+1} + Q E_{n,n+1}$ shifts the last column of $\mathrm{diag}(g,1)$ by $Q g_n$ (the $n$-th column of $g$ scaled by $Q$), which exactly compensates for the non-trivial central character twist introduced by the conductor.

Following Matringe (2012) and Jacquet (2009), one constructs $W_0 \in \mathcal{W}(\Pi,\psi^{-1})$ supported on $N_{n+1} A_{n+1} K_{n+1}(\mathfrak{q})$ (an Iwahori-type open compact subgroup of $G'$ depending on $\mathfrak{q}$) with the property that
\[
  \Psi(s;\,W_0,\,V_0) = L(s,\Pi\times\pi).
\]
The key point is that the $u_Q$-twist of the argument of $W$ makes the integrand $K_n(\mathfrak{q})$-equivariant, so that only the $V_0$-component (newvector) contributes.  One verifies: (a) convergence for all $s$ because the integral reduces to a finite sum over double cosets; (b) the value equals $L(s,\Pi\times\pi)$ by the JPSS functional equation for the $L$-factor together with the fact that the $\varepsilon$-factor arising from the $u_Q$ twist is $\varepsilon(s,\pi,\psi)\cdot\varepsilon(s,\Pi,\psi^{-1})^{-1}$, which cancels with the ratio of $L$-factors to give $L(s,\Pi\times\pi)$ exactly when both are evaluated at the new-/essential-vector level.
\end{proof}

\subsection*{Conclusion}

Since $L(s,\Pi\times\pi)$ is a rational function of $q^{-s}$ that is \emph{entire} (i.e., has no poles in $s$) and \emph{nonzero for all $s \in \mathbb{C}$} (by the general theory of local $L$-factors for generic representations: see JPSS 1983, Theorem 2.7, for the statement that $L(s,\Pi\times\pi)$ equals a product of factors of the form $(1 - \alpha q^{-s})^{-1}$ with $\alpha \in \mathbb{C}$, hence is never zero), we conclude:

Taking $W = W_0$ from Theorem~\ref{thm:essential} and $V = V_0$ (the newvector of $\pi$), the integral $\Psi(s;W,V) = L(s,\Pi\times\pi)$ is finite (the integral converges for all $s$ once written as a finite sum over cosets) and nonzero for all $s \in \mathbb{C}$.\hfill$\square$

\begin{remark}
The condition that $V \in \mathcal{W}(\pi,\psi)$ is unspecified in the problem (``for some $V$''), so the existence of $W$ alone is what is required.  The choice $V = V_0$ (newvector) and $W = W_0$ (essential vector) achieves this.
\end{remark}

\begin{remark}
The nontriviality of the twist $u_Q$ (vs.\ the untwisted integral with $u_Q = I$) is important when $\Pi$ is ramified.  For unramified $\Pi$ and $\pi$, the twist is trivial and the result reduces to the classical spherical case.
\end{remark}

\end{document}
